Un condensateur de capacité C $= 15\times10^{-5}$ Farads, initialement chargé à une tension de 10 Volts, se décharge à partir de l'instant $t_0=0$ à travers un circuit de résistance R$ = 2\times10^4$ Ohms. 

\smallskip

On sait que la tension $u$ est une fonction de la variable réelle $t$, exprimée en seconde, solution de l'équation différentielle suivante :

$$(E) : \qquad R~C~\dfrac{\text{d}u}{\text{d}t}(t)+u(t)=0$$



\begin{enumerate}
	\item Justifier que (E) peut s'écrire : $3u'+u=0$\par \vspace*{0.5cm}

	\Lignes{3}
	\item Résoudre l'équation homogène (E) {\it \small (On remarquera que c'est une équation homogène)}.\par \vspace*{0.5cm}

	\Lignes{2}
%	\item A l’aide du logiciel Géogébra, déterminer les nombres réels $a$ et $b$ pour que la fonction $p$ définie sur $\R$ par : $p(x) = a + b~x$ soit une solution particulière de l'équation différentielle (E), et donner l’expression de $p$. 
%	
%	{\it On prendra des curseurs allant de $0$ à $40$ par pas de $2$}\\
%	
%	On trouve : $a=$\makebox[0.1\linewidth]{\dotfill} \hfill $b=$\makebox[0.1\linewidth]{\dotfill} \hfill $p(x)=$\makebox[0.3\linewidth]{\dotfill}
%
%	\smallskip
%	
%	\hfill \fbox{ \bf Appeler le professeur}
	
	\item Justifier que $u(0)=10$\par \vspace*{0.5cm}

	\Lignes{2}
	\item Déterminer, par un calcul, la solution $u$ de (E) telle que $u(0)=10$.\par \vspace*{0.5cm}

	\Lignes{3}
	
	\item À l'aide de Géogebra, déterminer à partir de quel instant, que l'on notera $t_1$, $u()\leq \dfrac{10}{100}u_0$ ?
	
	\smallskip
	
	$t_1=$\makebox[0.1\linewidth]{\dotfill}

	\smallskip
	
	\appelprof
	\item On rappelle que la valeur moyenne d'une fonction $f$ entre $a$ et $b$, vaut $\ds{\dfrac{1}{b-a}\int_a^bf(x)\text{d}x}$
	\begin{enumerate}
		\item Rappeler la formule permettant de calculer $\ds{\int_a^bf(x)\text{d}x}$ à l'aide d'une primitive F de $f$. \par \vspace*{0.5cm}

	\Lignes{1}
		
		\item Déterminer une primitive U de $u$.\par \vspace*{0.5cm}

	\Lignes{3}
		\item À l'aide de Geogebra, donner une valeur approchée, au dixième de Volt près, de la valeur moyenne de la tension $u$ entre les instants $t_0$ et $t_1$.\par \vspace*{0.5cm}

	\Lignes{1}
		
	\end{enumerate}
\end{enumerate}